\chapter{Mongoose Schemas and Models}

\section{Schema vs Model}

\whatwhyhow{
A \textbf{schema} is the blueprint describing a document's shape, field types, and validation rules. A \textbf{model} is the compiled, usable interface Mongoose builds from that schema to actually query a collection.
}{
Separating the two lets you define validation and structure once (schema) while getting a rich query API (model) for free: \texttt{find}, \texttt{create}, \texttt{findById}, \texttt{updateOne}, and so on.
}{
Call \texttt{new mongoose.Schema(\{...\})} to define fields, then \texttt{mongoose.model("Task", taskSchema)} to compile it into a model bound to the \texttt{tasks} collection (Mongoose pluralizes and lowercases the name automatically).
}

\section{Task Schema}

\begin{lstlisting}[language=JavaScript]
// models/Task.js
const mongoose = require("mongoose");

const taskSchema = new mongoose.Schema(
  {
    title: {
      type: String,
      required: [true, "Title is required"],
      trim: true,
    },
    description: {
      type: String,
      default: "",
    },
    status: {
      type: String,
      enum: ["todo", "progress", "completed"],
      default: "todo",
    },
    priority: {
      type: String,
      enum: ["low", "medium", "high"],
      default: "medium",
    },
    dueDate: {
      type: Date,
    },
    user: {
      type: mongoose.Schema.Types.ObjectId,
      ref: "User",
      required: true,
    },
  },
  { timestamps: true }
);

const Task = mongoose.model("Task", taskSchema);

module.exports = Task;
\end{lstlisting}

\section{Field Options Explained}

\begin{table}[h]
\centering
\begin{tabularx}{\textwidth}{l X}
\toprule
\textbf{Option} & \textbf{Meaning} \\
\midrule
\texttt{type}     & The data type Mongoose casts and validates against (\texttt{String}, \texttt{Number}, \texttt{Date}, \texttt{Boolean}, \texttt{ObjectId}, ...). \\
\texttt{required}  & Rejects document creation/save if the field is missing; can pair with a custom error message. \\
\texttt{enum}      & Restricts a String/Number field to a fixed set of allowed values. \\
\texttt{default}   & Value used when the field is not supplied at creation time. \\
\texttt{ref}       & Names the model that an \texttt{ObjectId} field points to, enabling \texttt{.populate()}. \\
\texttt{trim}      & Strips leading/trailing whitespace from String values. \\
\texttt{timestamps} & Schema option (not a field option) that auto-adds and maintains \texttt{createdAt}/\texttt{updatedAt}. \\
\bottomrule
\end{tabularx}
\end{table}

\section{User Schema (for later chapters)}

\begin{lstlisting}[language=JavaScript]
// models/User.js
const mongoose = require("mongoose");

const userSchema = new mongoose.Schema(
  {
    name: {
      type: String,
      required: [true, "Name is required"],
      trim: true,
    },
    email: {
      type: String,
      required: [true, "Email is required"],
      unique: true,
      lowercase: true,
      trim: true,
    },
    password: {
      type: String,
      required: [true, "Password is required"],
      minlength: 6,
      select: false, // excluded from query results by default
    },
    avatar: {
      type: String,
      default: "",
    },
    role: {
      type: String,
      enum: ["user", "admin"],
      default: "user",
    },
  },
  { timestamps: true }
);

const User = mongoose.model("User", userSchema);

module.exports = User;
\end{lstlisting}

\begin{notebox}[NOTE]
\texttt{unique: true} on \texttt{email} creates a unique index at the database level; it is not a validation rule by itself and needs \texttt{unique} indexes to actually be built (Mongoose does this automatically on model compilation in development).
\end{notebox}

\begin{tipbox}[TIP]
\texttt{select: false} on \texttt{password} means \texttt{User.find()} never returns the password hash by default. You must explicitly opt in with \texttt{.select("+password")} when you need it, such as during login.
\end{tipbox}
